\section{Additional Related Works}

In addition to the related works on the NTK approach, mean-field analyses, and other works analyzing the training dynamics of neural networks that we discussed in \Cref{sec:intro}, we also discuss the following additional related works.

There are several works which study mean-field Langevin dynamics \citep{chen2023uniformintime, suzuki2023convergence}. These works show ``uniform-in-time propagation of chaos'' estimates, i.e. their results imply bounds on the coupling error between finite-width and infinite-width trajectories which do not grow exponentially in the time $T$. However, it is likely non-trivial to apply these analyses to directly obtain good test error and sample complexity in the setting of two-layer neural networks, since mean-field Langevin dynamics may require many iterations to obtain good population loss. More concretely, Theorem 4 of \citet{mei2018mean} suggests that the inverse temperature $\lambda$ for mean-field Langevin dynamics has to be at least proportional to the dimension $d$ in order for Langevin dynamics to achieve low population loss. This would cause the log-Sobolev constant in \citet{suzuki2023convergence} to be on the order of $e^{-d}$, leading to a running time of $e^{-d}$. In comparison, we show in \Cref{thm:projectedGD_main} that projected gradient descent with $\poly(d)$ iterations can achieve a low population loss. We also note that \citet{chen2023uniformintime} do not study the impact of finite samples on the coupling error.

We also note that there are several works which reduce the population dynamics to a lower-dimensional dynamics \citep{abbe2022merged, hajjar2023symmetries, arnaboldi2023highdimensional}, as we do in \Cref{subsec:pop_symmetry}. The focus of our work is on the analysis of the population dynamics (\Cref{sec:population_loss}) and coupling error between the empirical and population dynamics (\Cref{section:finite_sample_analysis}), rather than the dimension-free dynamics.

Finally, we note that our setting goes beyond that of \citet{abbe2022merged}: our target function does not satisfy the merged-staircase property required by \citet{abbe2022merged}, since the lowest-degree term in our target function has degree $2$. The work of \citet{abbe2023sgd} studies functions which have higher ``leap complexity,'' meaning that terms of the target function can introduce more than one new coordinate at once, or introduce a new coordinate using a Hermite polynomial with degree greater than $1$ (in the case where the inputs are Gaussian). However, \citet{abbe2023sgd} use a non-standard algorithm which separates the coordinates based on whether they are large or small, and applying different projections to each subset of coordinates.